% ===== 底本 第 3 章《环簇 Fano 集的 SYZ 镜像对称》；片段 tr04b；源块 B140--B154 =====

% [源 B140]
\begin{lemma*}
对任意 $\alpha\in\Omega^{\ast} (X_{\Sigma})$，求和
\[
  \sum_{\mu(\beta)=2k}\int_{\left[\mathcal{M}_1 (\beta)\right]^{vir}}\ev^{\ast}\alpha\cdot T^{\omega (\beta)}
\]
在 $k\geq 2$ 时为 $0$。
\end{lemma*}

% [Tutor Note: B140 是 B137「只有 Maslov 指标 $2$ 的圆盘影响曲率」的技术依据。积分是对模空间的虚拟基本类 $[\mathcal{M}_1 (\beta)]^{vir}$ 做的，被积式是 $\ev^{\ast}\alpha$（$\ev$ 把模空间映到 $X_{\Sigma}$）而不是 $\alpha$ 本身；$T^{\omega (\beta)}$ 里的 $\omega (\beta)$ 就是圆盘类 $\beta$ 的辛面积，与 B144--B145 的 $\omega (D(v_i))$ 同一记号，可据此核对两处没有用不同的面积约定。]

% [源 B141]
于是只有 Maslov 指标为 $2$ 的圆盘对曲率 $\mathfrak{m}^0$ 有非平凡贡献。注意它们对应于与某个环簇边界分支恰好相交一次的全纯圆盘，而事实上我们可以确定所有这样的全纯圆盘。

% [Tutor Note: B141 的「恰好相交一次」指相交重数为 $1$，而不是几何上只有一个交点：由 B135 与 B139，$\mu(u)=2\sum_i\mu_i$，故 $\mu(u)=2$ 当且仅当 $\sum_i\mu_i=1$，即只有一个 $\mu_i$ 等于 $1$、其余为 $0$。这条等价关系是核对 B142 分类完备性的直接依据。]

% [源 B142]
\begin{lemma*}
Maslov 指标为 $2$ 的全纯圆盘，全都是形如
\[
  \tilde{u} (z)=(c_1,\dotsc, c_{i-1}, c_i z, c_{i+1},\dotsc, c_N)
\]
的全纯圆盘 $\tilde{u}\colon (\mathbb{D}^2,\partial\mathbb{D}^2)\to (\mathbb{C}^N\setminus Z(\Sigma),\tilde{L})$ 的投影，其中 $c_j\in\mathbb{C}^{\ast}$，$1\leq i\leq N$。我们把这类圆盘的同伦类记作 $D(v_i)$。
\end{lemma*}

% [Tutor Note: B142 的「投影」指 $u=\pi\circ\tilde{u}$，其中 $\pi\colon \mathbb{C}^N\setminus Z(\Sigma)\to X_{\Sigma}$ 是自然映射（在稠密环面 $(\mathbb{C}^{\ast})^N$ 上即按 $T(\Sigma)$ 取商）。核对：$\tilde{u}(0)$ 只有第 $i$ 个坐标为零，故与 $D_i$ 交于一点、重数 $1$，与其余 $D_j$ 不交；由 B135 得 $\mu=2$，与 B139 的 $\sum_i\mu_i=1$ 一致。]

% [源 B143]
由于所有这些计数都以全纯圆盘的面积为权，我们还需要确定这些 Maslov 指标为 $2$ 的全纯圆盘的面积。

% [源 B144]
\begin{lemma*}
设 $b\in\mathbb{R}^n$ 是向量，则以 $L_b$ 为边界、同伦类为 $D(v_i)$ 的全纯圆盘的辛面积为
\[
  \omega (D(v_i))=2\pi (\langle b,v_i\rangle -\lambda_i),
\]
其中 $\lambda_i$ 是定义矩多面体时出现的常数。
\end{lemma*}

% [Tutor Note: B144 的 $\lambda_i$ 就是本章前面矩多面体定义 $P=\{x\in\mathbb{R}^n\vert \langle x,v_i\rangle\geq -\lambda_i\}$ 中的常数；特别地，环簇 Fano 取标准化矩多面体时所有 $\lambda_i=1$（见本章 Batyrev 判据一段）。核对方法：本式与 B145 第二个等号把 $T^{\omega (D(v_i))}$ 换成 $T^{2\pi(\langle b,v_i\rangle -\lambda_i)}$ 完全吻合，可据此确认 $2\pi$ 因子的位置——$T$ 的指数用的是辛面积 $\omega (D(v_i))$ 本身，系数 $2\pi$ 由面积单位带来。]

% [源 B145]
有了这些准备，我们就可以把 $\mathfrak{m}^0$ 显式写成如下形式：注意 $\mathfrak{m}^0$ 只统计只有一个输出（没有输入）的全纯圆盘，因此它在 $CF^{\ast} (L_b,\nabla;L_b,\nabla)$ 中代表一个元素，形式如下：
\begin{align*}
\mathfrak{m}^0 (L_b,\nabla)
&=\sum_{i=1}^N T^{\omega (D(v_i))}PD\left(\ev_{\ast}\left[\mathcal{M}_1 (D(v_i))\right]^{vir}\right)\hol_{\nabla} (\partial D(v_i))\\
&=\sum_{i=1}^N T^{2\pi(\langle b,v_i\rangle -\lambda_i)}\hol_{\nabla} (\partial D(v_i)),
\end{align*}
而函数 $z_i(L_b,\nabla)\colon =\exp\left(-\int_{D(v_i)}\omega\right)\hol_{\nabla} (\partial D(v_i))$（$1\leq i\leq N$）可以证明是全纯的。

% [Tutor Note: B145 第一个等号中，$PD\left(\ev_{\ast}\left[\mathcal{M}_1 (D(v_i))\right]^{vir}\right)$ 是把 $L_b$ 上的链经 Poincaré 对偶写成一个上链，$\hol_{\nabla} (\partial D(v_i))$ 则是平联络 $\nabla$ 沿圆盘边界闭路取的 holonomy；第二个等号就是把 B144 的面积公式代入（$T$ 是诺维科夫变量）。代入之后 $\mathfrak{m}^0$ 只依赖 $b$ 与 $\nabla$，即它是拉格朗日膜 $(L_b,\nabla)$ 的函数，这是下一步把它读成镜像环面上 Laurent 多项式的前提。]

% [Tutor Note: B145 只用了一个 $i$ 的求和，是因为 B142--B144 已经确定：Maslov 指标 $2$ 的圆盘恰有 $N$ 类 $D(v_1),\dots,D(v_N)$，每类对应一个边界分支 $D_i$。因此 $\mathfrak{m}^0$ 的 $N$ 项与除子/射线的 $N$ 个指标一一对应——这也解释了 B147 的 $W_{\Sigma}$ 为什么恰有 $N$ 个求和项（每根射线对应一个单项式）；注意镜像变量 $w=(w_1,\dots,w_n)$ 只有 $n$ 个，$N$ 是项数而非变量个数。]

% [Tutor Note: B145 末尾的 $z_i(L_b,\nabla)=\exp\left(-\int_{D(v_i)}\omega\right)\hol_{\nabla} (\partial D(v_i))$ 源文只说「可以证明是全纯的」，未给论证，此处照译；可核对的是：它只依赖 $b$ 与 $\nabla$，且 $\lvert z_i\rvert=\exp(-\omega (D(v_i)))$（对酉联络而言），这正是 B146 把它当作镜像 $(\mathbb{C}^{\ast})^n$ 坐标的依据。]

% [源 B146]
取定环面 $\left(\mathbb{C}^{\ast}\right)^n$ 的一组基 $\{w_i\}_{i=1}^w$，曲率项便在 $\left(\mathbb{C}^{\ast}\right)^n$ 上定义一个 Laurent 多项式，我们称之为 \emph{Hori-Vafa 超势}（Hori-Vafa superpotential）。同调镜像对称的另一部分则断言：

% [Tutor Note: B146 的 $\{w_i\}_{i=1}^w$ 中上限 $w$ 与基元素记号 $w_i$ 撞名，原文如此（已核对 raw HTML）。按「$(\mathbb{C}^{\ast})^n$ 的一组基」理解，上限应为 $n$（镜像环面维数），并非 B147 求和中的射线数 $N$；B147 的 $w_i^{v_i}$ 是同源的记号问题（标准写法为 $W_{\Sigma}=\sum_{i=1}^N w^{v_i}$，$w=(w_1,\dots,w_n)$ 为镜像坐标），两处均照录原文，解释以 B147 批注为准。又：Hori-Vafa 超势在此首次出现，按 glossary 译为「Hori-Vafa 超势」。]

% [源 B147]
\begin{statement*}
存在三角范畴的等价
\[
  D^{\pi}\mathcal{F} (X_{\Sigma},\omega_{\Sigma})\simeq HMF(\left(\mathbb{C}^{\ast}\right)^n,W_{\Sigma}),
\]
其中 $W_{\Sigma}\colon =\sum_{i=1}^N w_i^{v_i}$ 是 Hori-Vafa 超势。
\end{statement*}

% [Tutor Note: B147 的 $D^{\pi}\mathcal{F} (X_{\Sigma},\omega_{\Sigma})$ 是 $X_{\Sigma}$ 的导出 Fukaya 范畴（上标 $\pi$ 表示取分裂闭包，即 Karoubi 完备化），$HMF$ 是超势 $W_{\Sigma}$ 的矩阵分解范畴；这是同调镜像对称的 B 侧陈述，源文因此用无编号的「陈述」（statement*）而非定理。另：$w_i^{v_i}$ 中 $v_i$ 是扇的射线生成元（$\mathbb{Z}^n$ 中的向量），$w_i^{v_i}$ 表示相应单项式；该记号与 B146 把 $\{w_i\}$ 称作「$\left(\mathbb{C}^{\ast}\right)^n$ 的一组基」有混用（标准写法是 $W_{\Sigma}=\sum_{i=1}^N w^{v_i}$，$w=(w_1,\dots,w_n)$ 为镜像坐标），此处照录原文。]

% [源 B148]
回到 Fukaya $A_{\infty}$-代数 $CF^{\ast}(L_b,L_b)$ 的计算，我们想首先验证 Floer 上同调是良定义的，即 $\left(\mathfrak{m}^1\right)^2=0$。这可以由下面的命题实现：

% [源 B149]
\begin{proposition*}
假设 $L_b$ 不围出任何 Maslov 指标 $\leq 0$ 的全纯圆盘，则我们有：
\begin{enumerate}[leftmargin=1.9em,itemsep=1pt,topsep=2pt]
% [源 B150]
  \item $\mathfrak{m}^0$ 是点类 $[\mathrm{pt}]$ 的倍数；
% [源 B151]
  \item Floer 上同调 $HF^{\ast} (L,L)$ 良定义，且以 $[\mathrm{pt}]$ 作为一个上闭链（cocycle）；
% [源 B152]
  \item 链层面的乘积 $\mathfrak{m}^2$ 在 $HF(L,L)$ 上决定一个良定义的结合乘积，并且 $[\mathrm{pt}]$ 是该乘积的单位元。
  % [Tutor Note: B152 原文 "a well-defined associate product" 中 associate 当为 associative（结合的）之误（已核对 raw HTML 原文如此）；译文按「结合乘积」处理，照录原文未改。]
\end{enumerate}
\end{proposition*}

% [Tutor Note: B149 的假设「不围出 Maslov 指标 $\leq 0$ 的全纯圆盘」即 $L_b$ 的无障碍性条件：它排除了 Maslov 指标 $\leq 0$ 的圆盘带来的额外贡献，于是 $\mathfrak{m}^0$ 只剩 Maslov 指标 $2$ 的项（Maslov 指标 $2k$、$k\geq 2$ 的项由 B140 消失），故 B150 的「是点类倍数」才成立；源文没有展开这条推理。三条结论与 B153 的 $(\mathfrak{m}^1)^2=0$ 计算配套：第二条说 $[\mathrm{pt}]$ 是上闭链，第三条说 $[\mathrm{pt}]$ 是单位元。]

% [Tutor Note: B150 的「点类」$[\mathrm{pt}]$ 指 $HF^{\ast} (L_b,L_b)$ 中单位所在的类（de Rham 模型中即常数函数 $1$；链层面常记作点类）。因此「$\mathfrak{m}^0$ 是 $[\mathrm{pt}]$ 的倍数」即 $\mathfrak{m}^0=C\cdot[\mathrm{pt}]$ 对某个标量 $C=C(b,\nabla)$ 成立——B153 正是把这个标量记成 $C$。另：B151、B152 中的 $L$ 是 $L_b$ 的简写，原文如此，照译。]

% [源 B153]
特别地，我们有 $\left(\mathfrak{m}^1\right)^2=\mathfrak{m}^2 (\mathfrak{m}^0,x)+(-1)^{\deg (x)+1}\mathfrak{m}^2 (x,\mathfrak{m}^0)=((-1)^{\deg (x)}+(-1)^{\deg (x)+1})C=0$。

% [Tutor Note: B153 的第一个等号是 $A_\infty$ 关系在 $d=1$ 时的具体形式：把 $\mathfrak{m}^1\circ\mathfrak{m}^1$ 展开，得到含曲率 $\mathfrak{m}^0$ 的两项，符号 $(-1)^{\deg (x)+1}$ 是第二项的内符号。再代入 $\mathfrak{m}^0=C\cdot[\mathrm{pt}]$：末式 $((-1)^{\deg (x)}+(-1)^{\deg (x)+1})C=0$ 表明两项分别给出 $(-1)^{\deg (x)}C$ 与 $(-1)^{\deg (x)+1}C$，无论 $\deg (x)$ 取奇偶都恰相消——这正是用「$\mathfrak{m}^0$ 是点类倍数」推出 $\mathfrak{m}^1$ 为微分的关键一步。]

% [源 B154]
因此，对所有拉格朗日膜 $(L_b,\nabla)$，$HF^{\ast} ((L_b,\nabla),(L_b,\nabla))$ 都是良定义的。

% [Tutor Note: 末句用 $HF^{\ast} ((L_b,\nabla),(L_b,\nabla))$，而 B151 用 $HF^{\ast} (L,L)$：两者是同一对象，前者显式带上平联络 $\nabla$ 的数据。术语「拉格朗日膜」（Lagrangian brane）在此首现，指 $(L,\nabla)$ 这样的配对，按 glossary 译作「拉格朗日膜」。]
